\paragraph{Historical claim adjudication.}
Our unit of audit is a claim with a recorded decision prefix, an acting player, an allowed class of interventions, and a local endpoint predicate. For a feasibility claim, a certificate is an explicit sequence of engine-validated actions reaching that endpoint from the historical prefix. Such a witness establishes existence, not optimality or the behavior of a model-controlled continuation. An impossibility claim instead requires an upper bound covering every permitted inflow or an exhaustive search over the declared action class; unsuccessful search alone is insufficient.

For construction, we distinguish the submitted ordered operations from their target purchase multiset. Allocation-preserving repair may split count-valued operations into atomic purchases and reorder those purchases. A successful repair establishes that the target allocation was reachable, while the submitted action still failed the benchmark's intermediate-state legality requirements. We therefore report allocation feasibility and applied execution separately.

The present results are local adjudications of this finite corpus. They do not establish that exact-prefix auditing improves annotation accuracy on unseen environments. A direct evaluation can compare current-menu judgments, nominal resource bounds and executable reachability on a frozen set of recorded cases, retaining unknown outcomes and distinguishing positive witnesses from negative certificates. The manuscript build package regenerates exhibits from aggregate exports; reproducing behavioral adjudications additionally requires the saved prefixes, engine version and analysis environment.
